% Options for packages loaded elsewhere
\PassOptionsToPackage{unicode}{hyperref}
\PassOptionsToPackage{hyphens}{url}
\PassOptionsToPackage{dvipsnames,svgnames,x11names}{xcolor}
\documentclass[
  12pt,
  oneside,
  openany]{book}
\usepackage{xcolor}
\usepackage[a4paper,margin=2.5cm]{geometry}
\usepackage{amsmath,amssymb}
\usepackage{microtype}
\usepackage{graphicx}
\usepackage{booktabs}
\usepackage{enumitem}
\usepackage{titlesec}
\usepackage{fancyhdr}

% Page style: elegant running headers and clean footer (monochrome)
\pagestyle{fancy}
\fancyhf{}
\fancyhead[LE,RO]{\thepage}
\fancyhead[LO]{\nouppercase{\rightmark}}
\fancyhead[RE]{\nouppercase{\leftmark}}
\setlength{\headheight}{14pt}

\titleformat{\section}
  {\normalfont\bfseries\Large}
  {\thesection}{0.75em}{}

\titleformat{\subsection}
  {\normalfont\bfseries\large}
  {\thesubsection}{0.75em}{}

% Compact but readable lists
\setlist[itemize]{topsep=4pt,itemsep=2pt}
\setlist[enumerate]{topsep=4pt,itemsep=2pt}

% Slightly increased line spread for better readability
\linespread{1.1}
\setcounter{secnumdepth}{5}
\usepackage{iftex}
\ifPDFTeX
  \usepackage[T1]{fontenc}
  \usepackage[utf8]{inputenc}
  \usepackage{textcomp} % provide euro and other symbols
\else % if luatex or xetex
  \usepackage{unicode-math} % this also loads fontspec
  \defaultfontfeatures{Scale=MatchLowercase}
  \defaultfontfeatures[\rmfamily]{Ligatures=TeX,Scale=1}
\fi
\usepackage{lmodern}
\ifPDFTeX\else
  % xetex/luatex font selection
  \setmainfont[]{Cambria}
  \setmonofont[]{Consolas}
  \setmathfont[]{Cambria Math}
\fi
% Use upquote if available, for straight quotes in verbatim environments
\IfFileExists{upquote.sty}{\usepackage{upquote}}{}
\IfFileExists{microtype.sty}{% use microtype if available
  \usepackage[]{microtype}
  \UseMicrotypeSet[protrusion]{basicmath} % disable protrusion for tt fonts
}{}
\usepackage{setspace}
\makeatletter
\@ifundefined{KOMAClassName}{% if non-KOMA class
  \IfFileExists{parskip.sty}{%
    \usepackage{parskip}
  }{% else
    \setlength{\parindent}{0pt}
    \setlength{\parskip}{6pt plus 2pt minus 1pt}}
}{% if KOMA class
  \KOMAoptions{parskip=half}}
\makeatother
\usepackage{longtable,booktabs,array}
\newcounter{none} % for unnumbered tables
\usepackage{calc} % for calculating minipage widths
% Correct order of tables after \paragraph or \subparagraph
\usepackage{etoolbox}
\makeatletter
\patchcmd\longtable{\par}{\if@noskipsec\mbox{}\fi\par}{}{}
\makeatother
% Allow footnotes in longtable head/foot
\IfFileExists{footnotehyper.sty}{\usepackage{footnotehyper}}{\usepackage{footnote}}
\makesavenoteenv{longtable}
\pagestyle{fancy}
\setlength{\emergencystretch}{3em} % prevent overfull lines
\providecommand{\tightlist}{%
  \setlength{\itemsep}{0pt}\setlength{\parskip}{0pt}}
\usepackage{indentfirst}
\usepackage{microtype}
\usepackage{titlesec}
\usepackage{fancyhdr}
\titleformat{\section}[block]{\Large\bfseries}{\thesection.}{0.75em}{}
\titlespacing*{\section}{0pt}{3ex plus 1ex minus 0.2ex}{1.5ex plus 0.2ex}
\titleformat{\subsection}[block]{\large\bfseries}{\thesubsection.}{0.75em}{}
\titlespacing*{\subsection}{0pt}{2.5ex plus 1ex minus 0.2ex}{1.2ex plus 0.2ex}
\pagestyle{fancy}
\fancyhf{}
\cfoot{\thepage}
\renewcommand{\headrulewidth}{0pt}
\renewcommand{\chaptername}{}
\renewcommand{\chaptertitlename}{}
\titleformat{\chapter}[block]{\normalfont\huge\bfseries\raggedright}{\thechapter.}{0.5em}{}
\titlespacing*{\chapter}{0pt}{1.5em}{0.6em}
\usepackage{bookmark}
\IfFileExists{xurl.sty}{\usepackage{xurl}}{} % add URL line breaks if available
\urlstyle{same}
\hypersetup{
  pdftitle={Математический анализ},
  colorlinks=true,
  linkcolor={black},
  filecolor={Maroon},
  citecolor={Blue},
  urlcolor={black},
  pdfcreator={LaTeX via pandoc}}

\title{Математический анализ}
\usepackage{etoolbox}
\makeatletter
\providecommand{\subtitle}[1]{% add subtitle to \maketitle
  \apptocmd{\@title}{\par {\large #1 \par}}{}{}
}
\makeatother
\subtitle{Конспекты семинаров, 4-й семестр}
\author{}
\date{}

\begin{document}
\frontmatter
\maketitle

\renewcommand*\contentsname{Содержание}
{
\setcounter{tocdepth}{3}
\tableofcontents
\vspace{-1.25\baselineskip}
}
\setstretch{1.3}
% Avoid blank page: book's \mainmatter uses \cleardoublepage and the first
% \chapter also \clearpage — that inserts an extra empty sheet after the TOC.
\makeatletter
\renewcommand{\mainmatter}{%
  \@mainmattertrue
  \pagenumbering{arabic}}%
\makeatother
\mainmatter
\begin{center}\rule{0.5\linewidth}{0.5pt}\end{center}

\chapter{7-ой МАТ. Семинар
06.04.2026}\label{ux43eux439-ux43cux430ux442.-ux441ux435ux43cux438ux43dux430ux440-06.04.2026}

\section{Равномерная сходимость
(продолжение)}\label{ux440ux430ux432ux43dux43eux43cux435ux440ux43dux430ux44f-ux441ux445ux43eux434ux438ux43cux43eux441ux442ux44c-ux43fux440ux43eux434ux43eux43bux436ux435ux43dux438ux435}

\textbf{Определение равномерной сходимости:}
\[S_n(x) \overset{D}{\rightrightarrows} S(x)\]
\[\forall \varepsilon > 0 \ \exists N : \forall k \ge N \ \forall x_0 \in D \quad (|S_k(x_0) - S(x_0)| < \varepsilon)\]

Раскроем модуль в неравенстве:
\[-\varepsilon < S_k(x_0) - S(x_0) < \varepsilon\]
\[S(x_0) - \varepsilon < S_k(x_0) < S(x_0) + \varepsilon\]

Таким образом, определение можно переписать в следующем виде:
\[\underbrace{\forall \varepsilon > 0}_{\text{для любой $\varepsilon$-трубки}} \quad \underbrace{\exists N \ \forall k \ge N}_{\text{начиная с некоторого номера}} \quad \forall x_0 \in D \quad \underbrace{\left( S(x_0) - \varepsilon < S_k(x_0) < S(x_0) + \varepsilon \right)}_{(\ast)}\]

\textbf{Геометрический смысл \((\ast)\):} все графики \(y = S_k(x)\)
лежат в \(\varepsilon\)-трубке графика \(y = S(x)\).

\textbf{Что такое \(\varepsilon\)-трубка графика \(y = S(x)\) на
множестве \(D\)?} Это множество точек плоскости \(Oxy\), которые
заметаются графиком \(y = S(x)\) (при \(x \in D\)), когда он смещается
вверх и вниз на \(\varepsilon\). Границы трубки определяются кривыми
\(y = S(x) + \varepsilon\) и \(y = S(x) - \varepsilon\).

{\def\LTcaptype{none} % do not increment counter
\begin{longtable}[]{@{}
  >{\raggedright\arraybackslash}p{(\linewidth - 0\tabcolsep) * \real{0.0556}}@{}}
\toprule\noalign{}
\endhead
\bottomrule\noalign{}
\endlastfoot
\textbf{Из примера прошлого занятия:} Рассмотрим последовательность
\(S_k(x) = x^k\) на интервале \(x \in [0; 1)\). \emph{(По графику
видно)}, что при малых \(\varepsilon \in (0; 1)\) вообще все графики
\(y = S_k(x) = x^k\) выходят за пределы \(\varepsilon\)-трубки графика
\(y = S(x) \equiv 0\). При этом, если \(\varepsilon \ge 1\), то все
графики \(y = S_k(x)\) целиком лежат в \(\varepsilon\)-трубке
\(y = S(x)\). \\
\textbf{Вывод:} равномерной сходимости нет,
\(S_k(x) \underset{D}{\not\rightrightarrows} S(x)\). \\
*** \\
\#\#\# Пример 2 Дан функциональный ряд, частичная сумма (ч.с.) которого
равна: \(S_n(x) = \frac{\text{tg} x}{n}\) \\
\textbf{а) К чему поточечно сходится этот ряд на
\(D = \left(0; \frac{\pi}{2}\right)\)?} При фиксированном
\(x \in \left(0; \frac{\pi}{2}\right)\):
\(\lim_{n \to \infty} \frac{\text{tg} x}{n} = \lim_{n \to \infty} \frac{\text{tg} x}{\infty} = 0 \implies S(x) = 0 \quad \text{при } x \in \left(0; \frac{\pi}{2}\right)\) \\
\textbf{б) Есть ли равномерная сходимость на
\(D = \left(0; \frac{\pi}{2}\right)\)?} Так как предельная функция
\(S(x) = 0\) непрерывная, то, возможно, будет равномерная сходимость.
Проверим это. Рассмотрим графики первых частичных сумм:
\(S_1(x) = \text{tg} x\), \(S_2(x) = \frac{\text{tg} x}{2}\) и т.д. При
приближении \(x \to \frac{\pi}{2}\) графики уходят в бесконечность. \\
\textbf{Вывод:} Какую бы \(\varepsilon\)-трубку (вокруг \(y=0\)) мы ни
взяли, все \(S_k(x)\) уйдут за её границы в
\(\infty \implies S_k(x) \not\rightrightarrows S(x)\) (равномерной
сходимости нет). \\
\end{longtable}
}

\emph{(Примечание: далее разобран фрагмент нахождения \(N(\varepsilon)\)
для другой функции, судя по виду вычислений, для функции, содержащей
экспоненту).}

\textbf{3. б)} Графики \(S_n(x)\) располагаются целиком в
\(\varepsilon\)-трубке (начиная с некоторого номера)
\(\implies S_k(-1) < \varepsilon\)

\emph{Вычисления:} \[e^{-k+1} < \varepsilon\]
\[-k + 1 < \ln \varepsilon\] \[k > -\ln \varepsilon + 1\]
\[k > \ln \frac{1}{\varepsilon} + 1\]

Отсюда берем номер (округляя вверх):
\[N(\varepsilon) = \left\lceil \ln \frac{1}{\varepsilon} + 1 \right\rceil\]

\begin{center}\rule{0.5\linewidth}{0.5pt}\end{center}

\subsection{Пример 5}\label{ux43fux440ux438ux43cux435ux440-5}

Дана последовательность:
\[S_k(x) = x^k, \quad D = \left[0; \frac{1}{2}\right]\]

\textbf{а) Найти поточечную сходимость \(S_k(x) \xrightarrow{D} S(x)\)}
При \(x \in \left[0; \frac{1}{2}\right]\):
\[\lim_{k \to \infty} S_k(x) = \lim_{k \to \infty} x^k = 0 \quad (\text{т.к. } 0 \le x < 1)\]
\(\implies S(x) = 0\) при \(x \in \left[0; \frac{1}{2}\right]\).

\textbf{б) Исследовать на равномерную сходимость
\(S_k(x) \overset{D}{\rightrightarrows} S(x)\)} Да, равномерная
сходимость есть, т.к. \(\forall \varepsilon > 0\) график \(y = S_k(x)\)
будет целиком лежать в \(\varepsilon\)-трубке графика
\(y = S(x) \equiv 0\). Наибольшее значение на отрезке функция принимает
в правой точке: \[S_k\left(\frac{1}{2}\right) < \varepsilon\]
\[\Updownarrow\] \[\frac{1}{2^k} < \varepsilon\]
\[\frac{1}{\varepsilon} < 2^k\] \[k > \log_2 \frac{1}{\varepsilon}\]

\emph{(Пометка: просто взять
\(N(\varepsilon) = \left\lceil \log_2 \frac{1}{\varepsilon} \right\rceil\)
нельзя, т.к. при больших \(\varepsilon\) может получиться отрицательное
число).}

\textbf{в) Найти \(N(\varepsilon)\)} Для корректного задания номера
распишем его кусочно: * Если \(\varepsilon > \frac{1}{2}\), то берем
\(N(\varepsilon) = 1\) * Если \(\varepsilon \le \frac{1}{2}\), то берем
\(N(\varepsilon) = \lceil \log_2 \frac{1}{\varepsilon} \rceil + 10\)
\emph{(с запасом)}

Итак, искомый номер \(N(\varepsilon)\) задается системой:

\(N(\varepsilon) = 1\), если \(\varepsilon > \frac{1}{2}\)
\(N(\varepsilon) = \lceil \log_2 \frac{1}{\varepsilon} \rceil + 10\),
если \(\varepsilon \le \frac{1}{2}\)

\emph{(Скобки \(\lceil ... \rceil\) означают округление до целого числа
вверх).}

\backmatter
\end{document}
